\documentclass[10pt, russian]{beamer}
\usepackage[utf8]{inputenc}
\usepackage[russian]{babel}
\usepackage[T2A]{fontenc}
\usepackage{graphicx} % для картинок
\usepackage{booktabs} % для красивых таблиц

\usetheme{Madrid}
\usecolortheme{default}

\title{История Орловского государственного университета имени И.С. Тургенева}
\author{М. А. Семенов}
\institute{Орловский государственный университет имени И.С. Тургенева}
\date{18 мая 2026}

\begin{document}

% ===================== Титульный слайд =====================
\begin{frame}
    \titlepage
\end{frame}

% ===================== Слайд 1 =====================
\begin{frame}{Возникновение университета}
    \begin{itemize}
        \item Идея создания высшего учебного заведения в Орле появилась ещё в годы Первой мировой войны.
        \item 23 июня \textbf{1919} года — официальная дата рождения университета (Орловский пролетарский университет).
        \item Первый ректор — Николай Иосифович Конрад.
        \item В 1920 году на базе пролетарского университета и Института народного образования создан Орловский государственный университет.
    \end{itemize}
\end{frame}

% ===================== Слайд 2 =====================
\begin{frame}{Развитие в 1930-е годы}
    \begin{itemize}
        \item 5 августа 1931 года — создан Орловский индустриально-педагогический институт.
        \item Открытие института состоялось 16 октября 1931 года.
        \item В 1932 году физико-технический факультет преобразован в \textbf{физико-математический}.
        \item Университет активно развивался как педагогический вуз.
    \end{itemize}
\end{frame}

% ===================== Слайд 3 (с картинкой) =====================
\begin{frame}{Современный университет}
    \begin{columns}
        \column{0.55\textwidth}
            \begin{itemize}
                \item Крупнейший вуз Орловской области
                \item Более 19 000 студентов
                \item Многофакультетный классический университет
                \item Статус опорного университета региона
            \end{itemize}
        \column{0.45\textwidth}
            \centering
            \includegraphics[width=\textwidth]{oreluniver.jpg}
            \tiny{Главный корпус ОГУ}
    \end{columns}
\end{frame}

% ===================== Слайд 4 (таблица) =====================
\begin{frame}{Этапы развития университета}
    \begin{table}
        \centering
        \begin{tabular}{ll}
            \toprule
            \textbf{Год} & \textbf{Событие} \\
            \midrule
            1919 & Создание пролетарского университета \\
            1920 & Образование Орловского государственного университета \\
            1931 & Открытие индустриально-педагогического института \\
            1932 & Появление физико-математического факультета \\
            2015 & Реорганизация и присвоение имени И.С. Тургенева \\
            2016 & Получение статуса опорного университета \\
            \bottomrule
        \end{tabular}
    \end{table}
\end{frame}

% ===================== Слайд 5 =====================
\begin{frame}{Физико-математический факультет}
    \begin{itemize}
        \item Один из старейших факультетов университета.
        \item Готовит специалистов по математике, физике, информатике.
        \item В 2007 году на факультете обучалось около 700 студентов.
        \item Активно развивает научную и образовательную деятельность.
    \end{itemize}
\end{frame}

% ===================== Слайд 6 (заключительный) =====================
\begin{frame}{Заключение}
    \begin{center}
        \Large{Спасибо за внимание!}
        
        \vspace{1cm}
        \textbf{ОГУ имени И.С. Тургенева —} \\
        ведущий вуз Орловского региона с более чем 100-летней историей.
    \end{center}
\end{frame}

\end{document}